\setsubtitle{author : tobiichi3227}
\section{D. Merge, Rebase, Squash!}

\begin{frame}{\btitle{題目敘述}}
    \begin{itemize}
        \item \texttt{v2.0} 上有 $H$ 個 hunk；第 $h$ 個內容是 $U_h$，解 conflict 的成本是 $w_h$。
        \item \texttt{feat/fancy} 有 $N$ 個 commits；第 $i$ 個把 hunk $x_i$ 改成 $v_i$。
        \item 在相鄰 commits 間放 $K-1$ 個切點，得到恰好 $K$ 個 squash commits，再依序 rebase。
    \end{itemize}
\end{frame}

\begin{frame}{\btitle{子任務}}
    \begin{enumerate}
        \item{$N,\ H\le 15$}
        \item{$K = 1$}
        \item{$K = N$}
        \item{每個 hunk 最多修改一次}
        \item{$N \le 1000$}
        \item{無額外限制}
    \end{enumerate}
\end{frame}

\begin{frame}{\btitle{Conflict}}
    一個 patch $\text{old}\to\text{new}$ 遇到目前內容 $\text{current}$ 時：
    \[
        \text{conflict}
        \iff
        \text{current}\notin\{\text{old},\text{new}\}.
    \]

    目標：選擇切點，使 conflict 的總成本最小。
\end{frame}

\begin{frame}{\btitle{記號}}
    \begin{itemize}
        \item $S_i[h]$：執行完前 $i$ 個 commits 後，\texttt{feat/fancy} 上 hunk $h$ 的內容。
        \item $f_h$：hunk $h$ 第一次被修改的 commit；若從未修改則不存在。
    \end{itemize}
\end{frame}

\begin{frame}{\btitle{Squash 後的 patch}}
    區間 $(j,i]$ squash 後，hunk $h$ 的 patch 是
    \[
        S_j[h]\longrightarrow S_i[h].
    \]

    題目保證 hunk 不會被改回 $0$，所以包含 $f_h$ 的 squash commit 一定真的會修改 $h$。
\end{frame}

\begin{frame}{\btitle{第一次修改}}
    考慮包含 $f_h$ 的第一個 squash commit，假設其右端點為 $i$：
    \[
        0\longrightarrow S_i[h].
    \]

    此時 \texttt{v2.0} 的內容仍是 $U_h$。又因 $U_h\neq 0$，所以
    \[
        h\text{ 發生 conflict}
        \iff S_i[h]\neq U_h.
    \]
\end{frame}

\begin{frame}{\btitle{不會再次 conflict}}
    解完後目前內容變成 $S_i[h]$。之後每個 squash commit 的
    \texttt{old} 都等於目前內容，因此不可能再次 conflict。

    上面很抽象，可以感性理解一下，顯然你 Merge 時只要在某個 commit 解決衝突後，後面的 commit 就不可能衝突了。

    因為後面都是你自己的 commit，你不會 commit 會衝突的東西上去，這樣要怎麼 commit 成功。

    \begin{alertblock}{結論}
        hunk $h$ 的成本，只由「包含第一次修改 $f_h$ 的那一段」決定。
    \end{alertblock}
\end{frame}

\begin{frame}{\btitle{一段 squash commit 的成本}}
    定義 $C(j+1,i)$ 為把 commits $j+1,\ldots,i$ squash 成一段時，
    這一段新產生的 conflict 成本。

    它只需要計算第一次修改落在這段內的 hunks：
    \[
        \boxed{
        C(j+1,i)=
        \sum_{\substack{h:\;j<f_h\le i\\S_i[h]\neq U_h}} w_h
        }
    \]

    因為每個 hunk 只會在包含 $f_h$ 的那段被計算一次，
    所以所有段的 $C$ 相加，恰好就是整次 rebase 的總成本。
\end{frame}

\begin{frame}{\btitle{$N,\ H \le 15$}}
    從 $N-1$ 個空隙選出 $K-1$ 個切點：
    \[
        \binom{N-1}{K-1}\text{ 種分法}.
    \]

    對每種分法直接模擬：
    \begin{enumerate}
        \item 預先算出每個 commit 後的狀態 $S_i$。
        \item 對每段 squash commit 產生 patch。
        \item 按題目的規則 rebase，累加 conflict 成本。
    \end{enumerate}

    時間複雜度
    \[
        O\!\left(\binom{N-1}{K-1}(N+KH)\right),
    \]
    在 $N,\ H\le 15$ 時足夠。
\end{frame}

\begin{frame}{\btitle{$K = 1$}}
    全部 squash 成 $(0,N]$，答案為
    \[
        \sum_{h:\;f_h\text{ 存在},\ S_N[h]\neq U_h}w_h
    \]

    複雜度 $O(N+H)$。
\end{frame}

\begin{frame}{\btitle{$K = N$}}
    每個 commit 單獨套用，只有第一次修改可能 conflict。答案為
    \[
        \sum_{h:\;f_h\text{ 存在},\ v_{f_h}\neq U_h}w_h.
    \]

    複雜度 $O(N+H)$。
\end{frame}


\begin{frame}{\btitle{每個 hunk 最多修改一次}}
    不論怎麼切，包含那次修改的區間終點值都相同；答案同樣是
    $\sum_{h:\;f_h\text{ 存在},\ v_{f_h}\neq U_h}w_h$。

    複雜度 $O(N+H)$。
\end{frame}

\begin{frame}{\btitle{$N \le 1000$}}
    定義
    \[
        dp[k][i]
        =\text{前}\ i\text{ 個 commits 分成恰好}\ k\text{ 個非空段的最小成本}.
    \]

    初始條件：
    \[
        dp[0][0]=0,\qquad dp[0][i]=+\infty\quad(i>0).
    \]
\end{frame}

\begin{frame}{\btitle{$N \le 1000$}}
    枚舉最後一段前的切點 $j$：
    \[
        dp[k][i]
        =\min_{k-1\le j<i}
        \bigl(dp[k-1][j]+C(j+1,i)\bigr).
    \]

    最終答案是 $dp[K][N]$。
\end{frame}

\begin{frame}{\btitle{$N \le 1000$}}
    固定右端點 $i$，把目前滿足 $S_i[h]\neq U_h$ 的 hunk，
    依照第一次修改位置 $f_h$ 分桶：
    \[
        B_i[t]=\sum_{h:\;f_h=t,\ S_i[h]\neq U_h}w_h.
    \]

    每前進一個 commit，只會改變一個 hunk 所在桶子的值。
\end{frame}

\begin{frame}{\btitle{$N \le 1000$}}
    由成本公式可得
    \[
        C(l,i)=\sum_{t=l}^{i} B_i[t].
    \]

    再對 $B_i$ 做一次後綴和，就能在 $O(N)$ 算出所有 $C(l,i)$。

    \begin{itemize}
        \item 預處理：$O(N^2+H)$
        \item DP：$O(KN^2)$
    \end{itemize}
\end{frame}

\begin{frame}{\btitle{無額外限制}}
    固定 DP 層 $k$，令
    \[
        prev[j]=dp[k-1][j],
        \qquad
        A_i[j]=prev[j]+C(j+1,i).
    \]

    則
    \[
        dp[k][i]=\min_{k-1\le j<i}A_i[j].
    \]
\end{frame}

\begin{frame}{\btitle{無額外限制}}
    從左到右掃描右端點 $i$。我們需要維護所有切點 $j$ 的 $A_i[j]$，
    並支援：
    \begin{itemize}
        \item 一段切點同時加上相同成本；
        \item 查詢合法切點 $[k-1,i-1]$ 的最小值。
    \end{itemize}

    這可以用線段樹維護。
\end{frame}

\begin{frame}{\btitle{無額外限制}}
    commit $i$ 修改 hunk $h$，令 $f=f_h$。

    hunk $h$ 會被 $C(j+1,i)$ 計入的必要條件是
    \[
        j<f.
    \]

    因此因 $S_i[h]\neq U_h$ 造成改變時，
    它只會同時影響$j=0,1,\ldots,f-1$，這恰好是一段前綴。

    預先替每個 commit 算出
    \[
        right[i]=f_h-1,\qquad delta[i]=\text{成本變化},
    \]
    掃到 $i$ 時對 $[0,right[i]]$ 加上 $delta[i]$。
\end{frame}

\begin{frame}{\btitle{無額外限制}}
    第一次修改 hunk $h$（$i=f_h$）：
    \[
        right[i]=i-1,
    \]
    \[
        delta[i]=
        \begin{cases}
            +w_h,&v_i\neq U_h,\\
            0,&v_i=U_h.
        \end{cases}
    \]
\end{frame}

\begin{frame}{\btitle{無額外限制}}
    不是第一次修改，令 $old=S_{i-1}[h]$、$new=S_i[h]$：
    \[
        right[i]=f_h-1,
    \]
    \[
        delta[i]=
        \begin{cases}
            +w_h,&old=U_h,\ new\neq U_h,\\
            -w_h,&old\neq U_h,\ new=U_h,\\
            0,&\text{其餘情況}.
        \end{cases}
    \]

    也就是只追蹤「目前內容是否不等於 $U_h$」的狀態變化。
\end{frame}

\begin{frame}{\btitle{無額外限制}}
    對 $k=1,2,\ldots,K$：
    \begin{enumerate}
        \item 用 $prev[0],\ldots,prev[N-1]$ 建立線段樹。
        \item 依序掃描 $i=1,2,\ldots,N$。
        \item 對 $[0,right[i]]$ 加上 $delta[i]$。
    \end{enumerate}
\end{frame}

\begin{frame}{\btitle{無額外限制}}
    \begin{enumerate}
        \setcounter{enumi}{3}
        \item 若 $i\ge k$，查詢
        \[
            cur[i]=\min_{k-1\le j<i}A_i[j]
                  =\operatorname{rangeMin}(k-1,i-1).
        \]
        \item 掃完後交換 $prev$ 與 $cur$。
    \end{enumerate}

    初始只有 $prev[0]=0$，其餘皆為 $+\infty$；答案是 $prev[N]$。
\end{frame}

\begin{frame}{\btitle{無額外限制}}
    \begin{block}{線段樹不變量}
        在某一層掃描完 commit $i$ 後，葉節點 $j$ 儲存的值恰為
        \[
            prev[j]+C(j+1,i).
        \]
    \end{block}

    每個 hunk 對 $C(j+1,i)$ 的貢獻，只在它的狀態跨過 $U_h$ 時改變，
    且只影響 $j<f_h$；因此 prefix range add 完整維護了不變量。

    查詢 $[k-1,i-1]$ 就恰好枚舉所有合法的最後切點。
\end{frame}

\begin{frame}{\btitle{無額外限制}}
    時間複雜度 $O(H+KN\log_2 N)$，空間複雜度 $O(H+N)$。

    成本總和最大可達 $2\times10^{14}$，要用 long long。
\end{frame}